% ============================================================
%  致  谢
% ============================================================
\chapter*{致\ \ 谢}
\addcontentsline{toc}{chapter}{致谢}

\zihao{-4}\songti\raggedright\setlength{\parindent}{2em}
时光荏苒，四年的本科生涯即将走到尾声，希望这篇论文能够为这段经历画上一个令人欣慰的句点。

在此，首先要向医工学院的付东辽老师致以最诚挚的感谢。从选题之初到方案落地，付老师始终给予我充分的信任与支持，这份信任是我能够将设计推进到底的基石。课题的完成与论文的撰写，亦离不开付老师在此期间的悉心指导。

其次，感谢国际教育学院/莫动理工学院科创实验室所提供的宝贵平台。本设计的绝大部分工作——从电路调试到论文写作——均是在实验室中完成的，但更重要的是，这里是我真正踏入电子设计与机械设计领域第一块实践土壤。在实验室度过的无数个日日夜夜里，一次次项目历练和竞赛经验的叠加，逐渐内化为我的工程能力，不仅支撑了毕业设计的完成，也为研究生阶段的学习奠定了坚实的基础。

再次，感谢在课题推进过程中所有给予过我帮助的人——我的室友、实验室的伙伴，以及将在后续把本设计带入生物医学工程创新设计大赛的王俊皓同学，他在软件设计方面做出了重要的贡献。同时，也要感谢我的家人，他们的支持与信任始终是我前行中不可或缺的力量。

本设计是一项没有现成参考原型、也几乎没有直接相关设计文献可供借鉴的创新性尝试，因此存在诸多不成熟和亟待改进之处。在接下来的时间里，我将继续对装置进行迭代与打磨，努力推动其朝向更为完善的形态演进，并争取尽早开展临床试验与治疗有效性评估，使本设计最终能够真正服务于脑卒中患者，成为一件实用且有价值的眼部康复设备。

最后，前路尚远，唯愿不负所学，不怠于行。


\clearpage

